% !PRINTED 262
\weberpage{262}
But if $n$ is a prime, it is enough to show that $\Phi_{1,0,n}$
contains no negative powers of $q$, since the same follows for
$\Phi_{n,0,1}$ by the interchange of $\omega$ with $\omega:n$.

The constant value which the function $\Phi_{a,c,\partial}$ takes is
determined from \gesperrt{one} term of the expansion.

In carrying out these calculations the expansions § 24, (11), are of use,
\weq{14}{\begin{aligned}
A &= q^{-\frac{n-1}{24}r} \prod_{1,\infty}^{h} \left(\frac{1+q^{n(2h-1)}}{1+q^{2h-1}}\right)^r \\
&\quad + q^{\frac{n-1}{24}r} \prod_{1,\infty}^{h} \left(\frac{1+q^{2h-1}}{1+q^{n(2h-1)}}\right)^r \\
B &= q^{-\frac{n+1}{24}s} \prod_{1,\infty}^{h} (1+q^{2h-1})^s (1+q^{n(2h-1)})^s \\
&\quad + \left(\frac{2}{n}\right)^{r+s} q^{\frac{n+1}{24}s} 2^s \prod_{1,\infty}^{h} \frac{1}{(1+q^{2h-1})^s (1+q^{n(2h-1)})^s} .
\end{aligned}}
The formulae do not always come out simplest when $r$, $s$ are
assumed as small as possible; rather it proves most expedient to add
the further condition
\weq{15}{\frac{(n-1)r}{12} \equiv \frac{(n+1)s}{12} \ (\mathrm{mod}\ 2)}
as well.

In this way we obtain, for the first seven odd primes, the
following determination of $A$ and $B$:
\[
\begin{aligned}
n &= 3, & A &= \left(\frac{u}{v}\right)^6 + \left(\frac{v}{u}\right)^6, & B &= (uv)^3 - \frac{8}{(uv)^3}, \\
n &= 5, & A &= \left(\frac{u}{v}\right)^3 + \left(\frac{v}{u}\right)^3, & B &= (uv)^2 - \frac{4}{(uv)^2}, \\
n &= 7, & A &= \left(\frac{u}{v}\right)^4 + \left(\frac{v}{u}\right)^4, & B &= (uv)^3 + \frac{8}{(uv)^3}, \\
n &= 11, & A &= \left(\frac{u}{v}\right)^6 + \left(\frac{v}{u}\right)^6, & B &= uv - \frac{2}{vu}, \\
n &= 13, & A &= \frac{u}{v} + \frac{v}{u}, & B &= (uv)^6 - \frac{64}{(uv)^6}, \\
n &= 17, & A &= \left(\frac{u}{v}\right)^3 + \left(\frac{v}{u}\right)^3, & B &= (uv)^4 + \frac{16}{(uv)^4}, \\
n &= 19, & A &= \left(\frac{u}{v}\right)^2 + \left(\frac{v}{u}\right)^2, & B &= (uv)^3 - \frac{8}{(uv)^3}.
\end{aligned}
\]
